% ---------------------------------------------------------------------------------------------
% Theory. Notation as in docs/notation_v2.md (proxy mean written \psi_i, residual \xi_i, to avoid a clash with the
% map f and the sampling fraction f_h). Full proofs: sections/appendix_proofs.tex. Bibliography: refs_theory.bib.
% The main estimator is the difference estimator (PATTERN-DIF); the plain Horvitz--Thompson estimator is psi = 0.
% ---------------------------------------------------------------------------------------------
\section{Theory}\label{sec:theory}

This section collects the statistical properties of the design and of its two estimators, the plain
Horvitz--Thompson plug-in and the model-assisted difference estimator. We work in the design-based
framework of survey sampling \citep{cochran1977,sarndal1992}: the pool $U$, the predictions
$\hat y_{mi}$, the probabilities $\hat p_{mi}$, the groups $G_i$ and the labels $Y_i$ are fixed
numbers, and the only randomness is the selection of the labelled sample. Two facts make this
framework convenient. First, the strata, the allocation and the proxy are computed from
$(\hat y_{mi},\hat p_{mi},G_i)_{i\in U}$ alone, so they are fixed before any label is seen.
Second, no model for $Y$ given the predictions is needed for validity; the proxy $\pi_i$ enters
the allocation and, for the difference estimator, the estimator itself, but only through label-free
quantities. Both estimators are unbiased for the totals for every proxy
(Section~\ref{sec:theory-asym}); a wrong proxy costs efficiency, not validity
(Sections~\ref{sec:theory-dif} and~\ref{sec:theory-proxy}). Section~\ref{sec:theory-asym} gives unbiasedness, consistency and
asymptotic normality; Section~\ref{sec:theory-structure} the structure of the influence values
that drives the gains; Section~\ref{sec:theory-dif} the comparison of the difference and the plain
estimator, including tuning of the difference coefficient and zero-variance strata;
Section~\ref{sec:theory-alloc} the allocation; Section~\ref{sec:theory-deff}
the design effect as a function of the disagreement rate. Proofs are in
Appendix~\ref{sec:theory-proofs}.

\subsection{Set-up}\label{sec:theory-setup}

Write $t_i(Y_i)=Y_iu_i+(1-Y_i)v_i\in\mathbb{R}^C$ for the vector of contributions of instance $i$
to the totals, $T=\sum_{i\in U}t_i(Y_i)$, and $\tau=T/N$. All entries of $u_i,v_i$ lie in $\{0,1\}$,
so $\|t_i\|_\infty\le 1$. Every comparison estimand is a smooth function of the totals,
$\theta=f(T)$, and is scale free in the sense that $f(T)=\phi(T/N)$: the group rates are ratios of
totals and $\theta^{ERR}_{mj}=(E_m-E_j)/N$. The function $\phi$ is infinitely differentiable on the
set where the denominators $P_g/N$ and $Q_g/N$ are positive. The \emph{influence value} of
instance $i$ for a scalar estimand $\theta=f(T)$ is
\begin{equation}\label{eq:influence}
z_i=\nabla f(T)^{\top}t_i(Y_i)=N^{-1}\nabla\phi(\tau)^{\top}t_i(Y_i).
\end{equation}
Strata $h=1,\dots,H$ partition $U$ into sets $U_h$ of sizes $N_h$, with $W_h=N_h/N$; $n_h$ labels
are drawn from $U_h$ by simple random sampling without replacement, independently across strata
(STSI). For any variable $x$ we write $\bar x_h=N_h^{-1}\sum_{U_h}x_i$,
$S^2_{x,h}=(N_h-1)^{-1}\sum_{U_h}(x_i-\bar x_h)^2$ and $f_h=n_h/N_h$, and $\bar x_{s_h}$, $s^2_{x,h}$
for the corresponding sample mean and sample variance, and $c_h=N_h^2(1-f_h)/n_h$. The plain
estimator and its variance estimator are
\begin{equation}\label{eq:ht}
\hat T=\sum_{h=1}^HN_h\bar t_{s_h},\qquad \hat\theta=f(\hat T),\qquad
\hat V=\sum_{h=1}^H N_h^2\,\frac{1-f_h}{n_h}\,s^2_{\hat z,h},\quad
\hat z_i=\nabla f(\hat T)^{\top}t_i(Y_i).
\end{equation}
$\hat T$ is the Horvitz--Thompson estimator \citep{horvitz1952} for the inclusion probabilities
$n_h/N_h$, and $\hat V$ is the usual linearisation (Taylor) variance estimator
\citep[Ch.~5]{sarndal1992}. Its population counterpart is the linearised variance
\begin{equation}\label{eq:linvar}
V(\theta;n)=\sum_{h=1}^HN_h^2\,\frac{1-f_h}{n_h}\,S^2_{z,h}=\sum_hc_hS^2_{z,h}.
\end{equation}

\paragraph{Proxy means and the difference estimator.} Let $\pi_i\in[0,1]$ be any proxy for $P(Y_i=1)$
that is a function of $(\hat y_{mi},\hat p_{mi},G_i)_{i\in U}$ only (the design uses
$\pi_i=M^{-1}\sum_m\hat p_{mi}$). Its \emph{proxy mean} is the expected contribution under
$Y_i\sim\mathrm{Bernoulli}(\pi_i)$,
\begin{equation}\label{eq:psi}
\psi_i=\pi_iu_i+(1-\pi_i)v_i,\qquad \xi_i=t_i(Y_i)-\psi_i=(Y_i-\pi_i)(u_i-v_i),
\end{equation}
with $\|\psi_i\|_\infty\le1$ and $\|\xi_i\|_\infty\le1$. Because every unit of the pool carries a
prediction, $\psi_i$ is known for all $i\in U$ and the proxy total $F=\sum_{i\in U}\psi_i$ is a
\emph{known constant}, not an estimate: the pool is the finite population of interest, so, unlike in
prediction-powered inference with an unlabelled sample from a superpopulation
\citep{angelopoulos2023ppi,angelopoulos2026ppipp}, there is no sampling error in the proxy total and no
$(1+n/N_u)$ factor. For any label-free vector $\psi_i$ (not necessarily a proxy mean) the
\emph{difference estimator} of the totals is
\begin{equation}\label{eq:dif}
\hat T_\psi=\sum_{i\in U}\psi_i+\sum_{h=1}^HN_h\big(\bar t_{s_h}-\bar\psi_{s_h}\big)
=\sum_{h=1}^H\big\{F_h+N_h\bar\xi_{s_h}\big\},\qquad F_h=\sum_{i\in U_h}\psi_i,\qquad
\hat\theta_\psi=f(\hat T_\psi).
\end{equation}
Componentwise, $\hat T_{\psi,c}=\sum_{i\in U}\psi_{ic}+\sum_hN_h\,\mathrm{mean}_{s_h}\{t_{ic}(Y_i)-\psi_{ic}\}$.
With $\psi\equiv0$ it is the plain estimator \eqref{eq:ht}; with the proxy mean \eqref{eq:psi} it is the
stratified difference estimator, the estimator of PATTERN-DIF (it is a model-assisted, generalised
difference estimator in the sense of \citealt{cassel1976} and \citealt[Ch.~6]{sarndal1992}, and the
stratified version of the difference estimator of prediction-powered inference). The residual influence value,
its sample version and the linearised variance are
\begin{align}
e_i&=\nabla f(T)^{\top}\xi_i=z_i-\tilde z_i,\qquad \tilde z_i=\nabla f(T)^{\top}\psi_i,\notag\\
\hat V_\psi&=\sum_hN_h^2\frac{1-f_h}{n_h}s^2_{\hat e,h},\qquad \hat e_i=\nabla f(\hat T_\psi)^{\top}\xi_i,\qquad
V_\psi(\theta;n)=\sum_hc_hS^2_{e,h}.\label{eq:resid}
\end{align}
Thus $z_i=\tilde z_i+e_i$: the plain influence value splits into the part $\tilde z_i$ that the proxy
predicts and the residual $e_i$ that the labels must supply. The plain estimator is the case $\tilde z\equiv0$,
and $V_0=V(\theta;n)$ in \eqref{eq:linvar}.

\subsection{Unbiasedness, consistency and asymptotic normality}\label{sec:theory-asym}

Asymptotic statements refer to a sequence of pools $U_\nu$, $\nu=1,2,\dots$, with $N_\nu\to\infty$
and a fixed number $H$ of strata; the index $\nu$ is suppressed. Let
$r^2=\sum_hW_h^2/n_h$, which is the order of $\operatorname{Var}(\hat\tau_c)$ for a generic total.
Throughout, $\psi$ is either $0$ or a proxy mean \eqref{eq:psi}, or any label-free vector with
$\|\psi_i\|_\infty\le1$.

\begin{assumption}\label{prop:assumptions}
\begin{enumerate}
\item[(A1)] \emph{Design.} Strata, sample sizes and the proxy $\pi$ (hence $\psi$) are functions of
$(\hat y_{mi},\hat p_{mi},G_i)_{i\in U}$ only, in particular not refitted on the labelled sample;
$2\le n_h\le N_h$; $\min_hn_h\to\infty$; and
$f_h\le 1-\varepsilon$ for some $\varepsilon>0$ and all strata that are not fully enumerated
(strata with $n_h=N_h$ are known exactly and contribute no variance).
\item[(A2)] \emph{Limits.} $\tau_\nu\to\tau_\infty$, and $P_g/N\ge c$, $Q_g/N\ge c$ for some
$c>0$, $g\in\{0,1\}$.
\item[(A3)] \emph{Non-degeneracy.} The linearised variance of the estimator in use,
$V_\psi(\theta;n)$ in \eqref{eq:resid} ($=V(\theta;n)$ for $\psi=0$), satisfies $V_\psi(\theta;n)\ge c\,r^2$ for some $c>0$.
\item[(A4)] \emph{Lindeberg-type condition.} For every stratum $h$, either
$n_h(1-f_h)S^2_{e,h}N^2\to\infty$, or the stratum's share of the variance,
$N_h^2(1-f_h)S^2_{e,h}/\{n_hV_\psi(\theta;n)\}$, tends to zero.
\end{enumerate}
\end{assumption}

(A1) excludes the use of labels in the design and in the proxy. (A2) asks that each group contains a non-vanishing
fraction of positives and negatives, so that the group rates are well defined. (A3) excludes
estimands whose influence values are (nearly) identically zero, for instance two models that
make identical predictions, in which case $\hat\theta=\theta=0$ trivially; for the difference estimator it
also excludes a proxy that is (nearly) exact, which can make $S^2_{e,h}$ vanish
(Remark~\ref{rem:zero-var}). (A4) is a simple
sufficient form of the condition of \citet{hajek1960} for the central limit theorem under simple
random sampling without replacement; because the contributions are bounded, Hájek's condition holds
in every stratum whose scaled variance diverges, and (A4) only rules out strata that carry a
non-vanishing share of the variance while containing a bounded effective number of observations.
Asymptotics for stratified sampling with a fixed or growing number of strata are treated by
\citet{bickel1984}; for the linearisation variance estimator see \citet{krewski1981}, whose regime
(many strata, few draws per stratum) is complementary to ours (few strata, many draws per stratum).

\begin{proposition}[Validity]\label{prop:validity}
Let $\hat T_\psi$, $\hat\theta_\psi$, $\hat V_\psi$ and $V_\psi(\theta;n)$ be as in \eqref{eq:dif}--\eqref{eq:resid}; $\psi\equiv0$
gives $\hat T,\hat\theta,\hat V,V(\theta;n)$ of \eqref{eq:ht}--\eqref{eq:linvar}.
\begin{enumerate}
\item[(a)] If $n_h\ge1$ for all $h$, then $E(\hat T_\psi)=T$ and
$\operatorname{Cov}(\hat T_\psi)=\sum_hN_h^2(1-f_h)n_h^{-1}\Sigma^\xi_h$, with $\Sigma^\xi_h$ the population
covariance matrix of $\xi_i=t_i(Y_i)-\psi_i$ in stratum $h$; if $n_h\ge2$, the sample covariance matrix
of $\xi_i$ gives an exactly unbiased estimator of $\operatorname{Cov}(\hat T_\psi)$. This holds for any strata
and any allocation that do not depend on the labels, and for any label-free $\psi$, in particular for
any proxy.
\item[(b)] Under (A1)--(A2), $\hat\theta_\psi-\theta=O_p(r)$; in particular $\hat\theta_\psi$ is consistent.
\item[(c)] Under (A1)--(A4), $(\hat\theta_\psi-\theta)/V_\psi(\theta;n)^{1/2}\to N(0,1)$ in distribution.
\item[(d)] Under (A1)--(A3), $\hat V_\psi/V_\psi(\theta;n)\to1$ in probability. Consequently, under
(A1)--(A4) the Wald interval $\hat\theta_\psi\pm z_{1-\alpha/2}\hat V_\psi^{1/2}$ has asymptotic coverage
$1-\alpha$.
\end{enumerate}
\end{proposition}

\emph{Proof sketch.} (a) is the classical unbiasedness of the stratified expansion estimator
\citep[Ch.~5]{cochran1977} applied to $\xi_i$, because $\hat T_\psi-T=\sum_hN_h(\bar\xi_{s_h}-\bar\xi_h)$ and the
known term $F$ cancels exactly. For (b), each coordinate of $\hat\tau_\psi-\tau$ has variance at most
$2r^2$ because the contributions are bounded; a second-order Taylor expansion of $\phi$ around
$\tau$, valid by (A2), gives $\hat\theta_\psi-\theta=\nabla\phi(\tau)^\top(\hat\tau_\psi-\tau)+O_p(r^2)$.
For (c), the leading term is a sum of $H$ independent stratum means; each is asymptotically
normal by Hájek's theorem, and (A3) makes the $O_p(r^2)$ remainder negligible relative to the
standard deviation. For (d), the stratum sample covariance matrices are consistent for $\Sigma^\xi_h$
and $\nabla\phi(\hat\tau_\psi)\to\nabla\phi(\tau_\infty)$; (A3) turns the resulting $o_p(r^2)$ error
into a ratio statement. $\square$

The estimator $\hat\theta_\psi$ of a ratio-type estimand is not exactly unbiased; its bias is $O(r^2)$, which by (A3) is of smaller
order than its standard deviation. The joint version for the vector of all comparison estimands
follows from the Cramér--Wold device, provided (A3) and (A4) hold for every fixed linear
combination. The proof uses only that the contributions are bounded and that the proxy is label-free;
it makes no use of the proxy being accurate, calibrated or even a probability.

\subsection{Structure of the influence values}\label{sec:theory-structure}

For a pair $(m,j)$ let $D_i=\hat y_{mi}-\hat y_{ji}\in\{-1,0,1\}$, and write
$\delta_g=TPR_{mg}-TPR_{jg}$ and $\gamma_g=FPR_{mg}-FPR_{jg}$ for the group-wise components, so
that $\theta^{TPR}_{mj}=\delta_1-\delta_0$ and $\theta^{FPR}_{mj}=\gamma_1-\gamma_0$.

\begin{proposition}[Influence values]\label{prop:influence}
For every instance $i$ with $G_i=g$,
\begin{equation}\label{eq:infl-forms}
z^{TPR}_i=s_g\,\frac{Y_i(D_i-\delta_g)}{P_g},\qquad
z^{FPR}_i=s_g\,\frac{(1-Y_i)(D_i-\gamma_g)}{Q_g},\qquad
z^{ERR}_i=\frac{D_i(1-2Y_i)}{N},
\end{equation}
with $s_1=1$, $s_0=-1$. The influence values of $\theta^{TPR}$ and $\theta^{FPR}$ sum to zero
within each group. For a stratum $h$ inside group $g$, let $\pi_h=\bar Y_h$,
$\rho_h=N_h^{-1}\sum_{U_h}|D_i|$ (the pair's disagreement rate in the stratum),
$r_h=N_h^{-1}\sum_{U_h}Y_i|D_i|$ and $c_h=N_h^{-1}\sum_{U_h}Y_iD_i$ (not to be confused with the factor $c_h$ of
\eqref{eq:linvar}, which is not used in this proposition). Then, with
$\kappa_h=N_h/(N_h-1)$,
\begin{align}
S^2_{z^{TPR},h}&=\frac{\kappa_h}{P_g^2}\Big\{(r_h-c_h^2)+\delta_g^2\pi_h(1-\pi_h)
-2\delta_g\,c_h(1-\pi_h)\Big\},\label{eq:stratvar-tpr}\\
S^2_{z^{ERR},h}&=\frac{\kappa_h}{N^2}\Big\{\rho_h-e_h^2\Big\},\qquad
e_h=N_h^{-1}\textstyle\sum_{U_h}D_i(1-2Y_i),\label{eq:stratvar-err}
\end{align}
and $S^2_{z^{FPR},h}$ is \eqref{eq:stratvar-tpr} with $Y$ replaced by $1-Y$, $P_g$ by $Q_g$ and
$\delta_g$ by $\gamma_g$. In particular, in a stratum in which $m$ and $j$ agree on every
instance (always the case for $K_i\in\{0,M\}$),
\begin{equation}\label{eq:agree-var}
S^2_{z^{TPR},h}=\kappa_h\Big(\frac{\delta_g}{P_g}\Big)^2\pi_h(1-\pi_h),\qquad
S^2_{z^{FPR},h}=\kappa_h\Big(\frac{\gamma_g}{Q_g}\Big)^2\pi_h(1-\pi_h),\qquad
S^2_{z^{ERR},h}=0 .
\end{equation}
\end{proposition}

The proof is a direct calculation (Appendix~\ref{sec:theory-proofs}). Equation \eqref{eq:agree-var}
is the structural fact behind the design. Where all models agree, the numerators of the
comparison estimands receive nothing, and the only remaining influence is the denominator term
$-s_g\delta_gY_i/P_g$: label variability $\pi_h(1-\pi_h)$ scaled by the square of the group-wise
target component. Comparisons between competitive models have $|\delta_g|$ of a few hundredths, so
these strata, which typically contain most of the pool, carry a variance that is smaller by two to three
orders of magnitude than the per-instance variance in the disagreement strata, where
\eqref{eq:stratvar-tpr} is dominated by $r_h$, the share of positives on which the pair disagrees.
For the error difference the agreement strata carry no variance at all.

The residual influence values of the difference estimator have an equally explicit form, and a
feature that the plain values lack: the label enters only through $Y_i-\pi_i$.

\begin{proposition}[Residual influence values]\label{prop:residual}
Let $d_i=\nabla f(T)^\top(u_i-v_i)$, which does not depend on the labels. Then $e_i=(Y_i-\pi_i)\,d_i$, and for $G_i=g$
\begin{align}
e^{TPR}_i&=s_g\frac{(Y_i-\pi_i)(D_i-\delta_g)}{P_g},\qquad
e^{FPR}_i=s_g\frac{\{(1-Y_i)-(1-\pi_i)\}(D_i-\gamma_g)}{Q_g},\notag\\
e^{ERR}_i&=-\frac{2(Y_i-\pi_i)D_i}{N},\label{eq:resid-forms}
\end{align}
and $\tilde z_i=z_i-e_i$ equals \eqref{eq:infl-forms} with $Y_i$ replaced by $\pi_i$. In a stratum where the pair agrees,
$e^{ERR}_i=0$ for every $i$: the stratum contributes neither variance nor error to the difference estimator of $\theta^{ERR}$.
Under the working model that $Y_1,\dots,Y_N$ are independent with $P(Y_i=1)=\pi_i$ (the proxy is calibrated), with expectation $E_\pi$,
the gradient evaluated at the non-random model-expected totals $T_\pi=\sum_{i\in U}\psi_i$ (as in the design; this changes $d_i$ by a relative $O(r)$ only),
and with $\bar\sigma^2_h=N_h^{-1}\sum_{U_h}\pi_i(1-\pi_i)d_i^2$ and $C_h=\mathrm{Cov}_h(z,\tilde z)$ the stratum covariance,
\begin{equation}\label{eq:resid-moments}
E_\pi(e_i)=0,\quad \mathrm{Var}_\pi(e_i)=\pi_i(1-\pi_i)d_i^2,\qquad
E_\pi S^2_{e,h}=\bar\sigma^2_h,\quad E_\pi S^2_{z,h}=S^2_{\tilde z,h}+\bar\sigma^2_h,\quad E_\pi C_h=S^2_{\tilde z,h}.
\end{equation}
\end{proposition}

\subsection{The difference estimator against the plain estimator}\label{sec:theory-dif}

This subsection answers four questions: when does $\hat\theta_\psi$ beat $\hat\theta$, by how much, what happens
if the proxy is wrong, and what do tuning of the coefficient and zero-variance strata do to inference. Write
$\beta_h=\mathrm{Cov}_h(z,\tilde z)/S^2_{\tilde z,h}$ for the within-stratum regression slope of the plain influence value on its proxy
mean, and $R_h$ for their within-stratum correlation, so that $\beta_h=R_hS_{z,h}/S_{\tilde z,h}$
(terms with $S^2_{\tilde z,h}=0$ are dropped). For a vector $\lambda=(\lambda_h)$ of label-free constants consider the family
$\hat T_{\lambda\psi}$ obtained by replacing $\psi$ by $\lambda_h\psi$ in stratum $h$; $\lambda\equiv0$ is the plain estimator and $\lambda\equiv1$
the difference estimator. By Proposition~\ref{prop:validity}(a) every member is unbiased for $T$.

\begin{proposition}[Variance of the difference estimator, design-based]\label{prop:dif-var}
Under the notation above, for any label-free proxy and any allocation:
\begin{enumerate}
\item[(a)] \emph{Identity.} $S^2_{e,h}=S^2_{z,h}-2\,\mathrm{Cov}_h(z,\tilde z)+S^2_{\tilde z,h}$, hence
\begin{equation}\label{eq:dif-identity}
V_0-V_\psi=\sum_hc_h\,S^2_{\tilde z,h}\,(2\beta_h-1)=\sum_hc_h\,S_{\tilde z,h}\,(2R_hS_{z,h}-S_{\tilde z,h}).
\end{equation}
The difference estimator has the smaller linearised variance if and only if the right side is positive. A sufficient condition is
$R_h>S_{\tilde z,h}/(2S_{z,h})$ in every stratum, that is $\beta_h>1/2$: the within-stratum correlation of $z_i$ with its proxy mean must exceed half
the ratio of their standard deviations.
\item[(b)] \emph{Coefficient.} $V(\lambda)=\sum_hc_h\{S^2_{z,h}(1-R_h^2)+(\lambda_h-\beta_h)^2S^2_{\tilde z,h}\}$. It is minimised at
$\lambda_h=\beta_h$, with $V(\beta)=\sum_hc_hS^2_{z,h}(1-R_h^2)$; moreover
$V_0=V(\beta)+\sum_hc_h\beta_h^2S^2_{\tilde z,h}$ and $V_\psi=V(\beta)+\sum_hc_h(1-\beta_h)^2S^2_{\tilde z,h}$. With the clipped coefficient
$\lambda_h=\min\{1,\max(0,\beta_h)\}$, $V(\lambda)\le\min(V_0,V_\psi)$.
\item[(c)] \emph{Calibrated proxy, design-based form.} If $\mathrm{Cov}_h(e,\tilde z)=0$ in every stratum (the residual is uncorrelated with the
proxy influence within strata), then $\beta_h=1$ and $V_0-V_\psi=\sum_hc_hS^2_{\tilde z,h}\ge0$ exactly.
\end{enumerate}
\end{proposition}

Part (a) shows that, unlike what is often assumed for prediction-powered estimators, the sign of the gain is not automatic:
$\lambda=1$ is the right coefficient only if the proxy influence is neither over- nor under-dispersed relative to
its regression on $z$; the harm from an over-dispersed proxy is quantified in Proposition~\ref{prop:dif-proxy}.
When the proxy is calibrated, however, the gain is guaranteed and can be computed without labels.

\begin{proposition}[Design-effect ratio under a calibrated proxy]\label{prop:dif-deff}
Under the working model of Proposition~\ref{prop:residual}, let $\nu_h=S^2_{\tilde z,h}/\bar\sigma^2_h$ (signal-to-noise ratio of the proxy
influence in stratum $h$; strata with $\bar\sigma^2_h=S^2_{\tilde z,h}=0$ are omitted from the sums, and $\bar\sigma^2_h>0$ is assumed otherwise) and $a_h=c_h\bar\sigma^2_h$. Then
\begin{enumerate}
\item[(a)] $E_\pi V_0-E_\pi V_\psi=\sum_hc_hS^2_{\tilde z,h}\ge0$, with equality if and only if $\tilde z$ is constant within every stratum that is not fully enumerated.
\item[(b)] For the same strata and sample sizes the design effect of the difference estimator relative to the plain estimator is
\begin{equation}\label{eq:dif-deff}
\frac{E_\pi V_\psi}{E_\pi V_0}=\frac{\sum_ha_h}{\sum_ha_h(1+\nu_h)}=\frac1{G},\qquad
G=\frac{\sum_ha_h\,(1+\nu_h)}{\sum_ha_h}=\frac{\sum_ha_h/(1-R_h^2)}{\sum_ha_h},
\end{equation}
where $R_h^2=\nu_h/(1+\nu_h)$ is the (expected) squared within-stratum correlation of $z_i$ with $\tilde z_i$. The gain $G$ is the
average of the stratum gains $1/(1-R_h^2)$ weighted by the strata's shares of the residual variance; $G=1+\nu$ if $\nu_h\equiv\nu$.
\item[(c)] For $\theta^{TPR},\theta^{FPR}$ and $\theta^{ERR}$ the stratum gains depend on the proxy as follows. Let $\overline{\pi(1-\pi)}$ denote
the stratum mean of $\pi_i(1-\pi_i)$. In a \emph{disagreement stratum} ($|D_i|=1$ on $U_h$, $D_i$ and $\pi_i$ uncorrelated within the stratum,
$\bar a$ the stratum mean of $D_i$, and $\delta_g,\gamma_g$ negligible),
\begin{align}
1+\nu^{TPR}&=\frac{\bar\pi(1-\bar\pi\bar a^2)}{\overline{\pi(1-\pi)}},\qquad
1+\nu^{FPR}=\frac{(1-\bar\pi)\{1-(1-\bar\pi)\bar a^2\}}{\overline{\pi(1-\pi)}},\notag\\
1+\nu^{ERR}&=\frac{1-\bar a^2(1-2\bar\pi)^2}{4\,\overline{\pi(1-\pi)}} .\label{eq:dif-closed}
\end{align}
In an \emph{agreement stratum} ($D_i=0$) the error difference has $e_i=\tilde z_i=0$, and for the rate differences
$1+\nu=\{\overline{\pi(1-\pi)}+S^2_{\pi,h}\}/\overline{\pi(1-\pi)}$, the inverse of the share of the Bernoulli variance that the proxy does \emph{not} explain.
\end{enumerate}
\end{proposition}

Four consequences. First, calibration alone guarantees $G\ge1$: with a calibrated proxy the difference estimator is never worse than the plain one in
expectation, for any allocation. Second, the gain is a within-stratum $R^2$, so it is governed by what the proxy influence $\tilde z_i$ explains of $z_i$
\emph{inside} a stratum, and the strata are already built to remove between-stratum variation. Third, $\tilde z_i$ contains two kinds of
information: the proxy $\pi_i$ and the known structure of $u_i,v_i$ (the signs and the pattern of $D_i$). The second kind is free and
is the main source of the gain: in a two-sided disagreement stratum with a constant proxy $\pi_i\equiv\bar\pi$ and $\bar a=0$,
\eqref{eq:dif-closed} gives gains $1/(1-\bar\pi)$ for $\theta^{TPR}$, $1/\bar\pi$ for $\theta^{FPR}$ and $1/\{4\bar\pi(1-\bar\pi)\}$ for $\theta^{ERR}$, i.e.\ 1.4,
3.3 and 1.2 for $\bar\pi=0.3$ and exactly 2, 2 and 1 for $\bar\pi=1/2$. Only the error difference needs a sharp proxy: its gain is
$\{1-\overline{(1-2\pi)^2}\}^{-1}$ and equals one for $\pi\equiv1/2$, which makes $\tilde z\equiv0$ and the two estimators identical. Fourth, the
gain is bounded for proxies of realistic sharpness, because the disagreement strata are the strata in which the compared models, and hence typically the
averaged probabilities $\pi_i$, are least certain: a large $G$ requires within-stratum $R_h^2$ close to one, that is a sharp proxy
($\pi_i$ near 0 or 1) and a stratum in which the models still disagree. Where the proxy is that sharp the predictions would be nearly
exact and few labels would be needed in any case. For disagreement strata with $\bar\pi$ between 0.2 and 0.5 the closed forms give gains of the order of one to three, on top of the
stratification gain; Section~\ref{sec:sim-dif} compares the prediction with the Monte Carlo ratios.

Gain and stratification multiply: for the same sample the efficiency of the difference estimator relative to SRS is, to the order of the
approximations of Section~\ref{sec:theory-deff}, $\mathrm{Eff}_t\cdot G_t$, with $\mathrm{Eff}_t$ from
\eqref{eq:deff-general} and $G_t$ from \eqref{eq:dif-deff}. They are not independent: a stratifier that
removes more of the proxy's within-stratum signal (a finer partition of the proxy range) reduces $G$ and raises $\mathrm{Eff}_t$.

\begin{proposition}[Misspecified proxy for the difference estimator]\label{prop:dif-proxy}
\begin{enumerate}
\item[(a)] For every proxy, $\hat T_\psi$ is design-unbiased and $\hat V_\psi$ is a consistent variance estimator
(Proposition~\ref{prop:validity}). Misspecification can affect efficiency only.
\item[(b)] Let the labels be independent with $P(Y_i=1)=q_i$, where $q_i$ may differ from $\pi_i$, and let $\tilde z^{q}_i$, $\tilde z^{\pi}_i$ be the proxy
influence values for $q$ and $\pi$. With $E_q$ the expectation under this model,
\begin{equation}\label{eq:dif-misspec}
E_q(V_0-V_\psi)=\sum_hc_h\Big\{S^2_{\tilde z^{q},h}-S^2_{\tilde z^{q}-\tilde z^{\pi},h}\Big\},
\end{equation}
so the difference estimator wins if and only if, within strata and on the $c_h$-weighted scale, the proxy error in influence units is smaller than the signal it
should capture. If $\tilde z^{\pi}_i-\bar{\tilde z}^{\pi}_h=\kappa(\tilde z^{q}_i-\bar{\tilde z}^{q}_h)$ in every stratum (a proxy that is sharper by the factor
$\kappa$ than the truth), the estimator wins if and only if $0<\kappa<2$; for $\kappa<0$ (anti-correlated proxy) it loses.
\item[(c)] In all cases $V_\psi\le(1+\max_hS_{\tilde z,h}/S_{z,h})^2\,V_0$, with equality only if $\tilde z_i=-s\,z_i$ within each stratum
for a constant $s>0$ (perfect anti-correlation). Hence the possible loss is bounded by a function of the dispersion of the proxy influence relative to the plain one.
\item[(d)] With the clipped coefficient of Proposition~\ref{prop:dif-var}(b) the variance never exceeds that of the plain estimator
(to first order). In particular, an uninformative proxy, in the sense $\beta_h=0$ ($\tilde z$ uncorrelated with $z$ within strata), has optimal coefficient $\lambda_h=0$ and
then $V_\psi=V_0+\sum_hc_hS^2_{\tilde z,h}$: the fixed coefficient $\lambda=1$ adds the proxy's own variance.
\end{enumerate}
\end{proposition}

Two clarifications. The statement that the difference estimator can be worse than the plain estimator \emph{only} if the proxy is anti-correlated is
too strong: by \eqref{eq:dif-identity} it is worse whenever $\beta_h<1/2$, so a weakly aligned or over-dispersed
proxy ($\kappa>2$) also hurts, while a modestly noisy but correctly scaled proxy does not. And ``uninformative'' must be read in terms of $\tilde z_i$,
not of $\pi_i$: a constant proxy $\pi_i\equiv\bar\pi$ still has $\tilde z_i$ varying with the known pattern of $D_i$, and for the rate differences
it still gains, as the closed forms \eqref{eq:dif-closed} show; only for $\theta^{ERR}$ with $\bar\pi=1/2$ is $\tilde z\equiv0$, and
only an inverted proxy ($\kappa<0$) makes $\beta_h<0$ for the error difference. Section~\ref{sec:sim-dif} confirms both points numerically.

\paragraph{Fixed versus tuned coefficient.} In prediction-powered inference \citep{angelopoulos2026ppipp,fisch2024stratppi} and in the
regression-estimator literature \citep{cassel1976,sarndal1992} the coefficient is estimated from the labelled sample, here
per stratum and per estimand: $\hat\lambda_h=\widehat{\mathrm{Cov}}(\hat z,\hat{\tilde z})/\widehat{\mathrm{Var}}(\hat{\tilde z})$ clipped to $[0,1]$.

\begin{proposition}[Estimated coefficient]\label{prop:dif-lambda}
Suppose $\hat\lambda_h$ is a function of the labelled sample of stratum $h$ with $\hat\lambda_h-\lambda^\infty_h\to0$ in probability, for deterministic
$\lambda^\infty_h$ (for the clipped sample slope, $\lambda^\infty_h=\min\{1,\max(0,\beta_h)\}$ whenever $\beta_h$ converges). Under (A1)--(A4)
for the fixed coefficient $\lambda^\infty$, $\hat\theta_{\hat\lambda\psi}-\hat\theta_{\lambda^\infty\psi}=o_p(V^{1/2})$; hence the tuned estimator has the same asymptotic
normal limit and linearised variance $V(\lambda^\infty)$ as the one with the optimal fixed coefficient, and is never asymptotically worse than the plain
estimator. At second order, estimating $\hat\lambda_h$ adds to the variance of stratum $h$ approximately
$c_h(1-R_h^2)S^2_{z,h}/n_h$, a relative increase of order $1/n_h$ in a stratum that is already optimally tuned, and the plug-in estimator $\hat V_\psi$ computed
with the fitted $\hat\lambda_h$ does not account for it (it is biased downwards by a factor of the same order).
\end{proposition}

The consequence is a trade-off that is independent of the particular data. The best that tuning can add to $\lambda=1$ is the factor
$V_\psi/V(\lambda^\ast)=1+\sum_hc_h(1-\beta_h)^2S^2_{\tilde z,h}/V(\beta)$ \eqref{eq:dif-identity}, which is small when the proxy is nearly calibrated ($\beta_h\approx1$): then
tuning gains little in variance but costs an $O(1/n_h)$ inflation in each of the $H$ strata and a downward-biased variance estimator, which matters when $n_h$ is small, as
at the budgets of this paper ($n_h$ of the order of 10). With a miscalibrated proxy the gain from tuning can be real (see the $\kappa\ne1$ cases in
Section~\ref{sec:sim-dif}), but it is paid for in coverage. The calibrated difference estimator with $\lambda=1$ has no estimated
parameters, exact unbiasedness and a plug-in variance that is unbiased for the totals (Proposition~\ref{prop:validity}(a)), which is why we adopt it as the main
estimator. Section~\ref{sec:sim-dif} reports the predicted and the observed gain of tuning together with the coverage of each version.

\begin{remark}[Zero-variance strata and the variance floor]\label{rem:zero-var}
The Wald interval needs $\hat V_\psi>0$ and $\hat V_\psi\approx V_\psi$, and both fail when a stratum's sample shows no variation in $\hat e_i$ although
$S^2_{e,h}>0$. For the difference estimator this has three sources. (i) \emph{Structural zeros}: strata with $d_i\equiv0$ (the error
difference where the pair agrees, Proposition~\ref{prop:residual}), where $e_i\equiv0$ in the whole stratum; the estimate is exact there and
$\hat V_\psi$ correctly contributes zero. (ii) \emph{Sampling zeros}: a stratum with few positives (or negatives) in which the sample has none. For the plain estimator $\hat z_i=0$ for every sampled unit, so $s^2_{\hat z,h}=0$ exactly;
for the difference estimator $\hat e_i=-\pi_i\hat d_i$ varies with $\pi_i\hat d_i$, so exact zeros are rarer, but the sample variance is still far below the
population variance $\bar\sigma^2_h=N_h^{-1}\sum_{U_h}\pi_i(1-\pi_i)d_i^2$ whenever the missed positives matter. With a calibrated proxy of level $\pi$ and constant $d$, the
probability of a sample without positives is $(1-\pi)^{n_h}$, which is the zero-count failure of the Wald interval for a binomial proportion \citep{brown2001}.
(iii) \emph{Exact proxies}: if $\pi_i\in\{0,1\}$ equals $Y_i$ then $e_i\equiv0$ and
(A3) fails; the estimator is then exact, but the proxy is unlikely to be exact on the whole pool. The floor of Section~\ref{sec:ablations} carries over: replace $s^2_{\hat e,h}$ by the anticipated
residual variance
\begin{equation}\label{eq:floor}
\tilde s^2_{h}=N_h^{-1}\sum_{i\in U_h}\pi_i(1-\pi_i)\hat d_i^{\,2},\qquad \hat d_i=\nabla f(\hat T_\psi)^\top(u_i-v_i),
\end{equation}
whenever $s^2_{\hat e,h}=0$ (or, more conservatively, use $\max\{s^2_{\hat e,h},\eta\tilde s^2_h\}$ with $\eta\in(0,1]$). Because $\tilde s^2_h=0$ exactly when $d_i\equiv0$ on $U_h$, structural zeros
are untouched. The floor is label-free, never decreases $\hat V_\psi$, leaves the point estimate unchanged, and (Appendix~\ref{sec:theory-proofs}) does not affect
Proposition~\ref{prop:validity}(c)--(d): the floor is active with probability tending to zero in every stratum with $S^2_{e,h}$ bounded away from zero, and in strata
with vanishing variance it is, if anything, conservative. It is a finite-sample safeguard, not an asymptotic correction, and Section~\ref{sec:sim-dif} reports its effect on coverage.
\end{remark}

\subsection{Allocation}\label{sec:theory-alloc}

Let $w_t\ge0$ be weights on a finite set of targets $t$ with linearised variances
$V_t(n)=\sum_hN_h^2(1/n_h-1/N_h)S^2_{ht}$, where $S^2_{ht}=S^2_{z,ht}$ for the plain estimator and $S^2_{ht}=S^2_{e,ht}$ for the difference estimator, and let
$A_h=(\sum_tw_tS^2_{ht})^{1/2}$. The weighted objective is
\begin{equation}\label{eq:objective}
\Phi(n)=\sum_tw_tV_t(n)=\sum_h\frac{N_h^2A_h^2}{n_h}-\sum_hN_hA_h^2 .
\end{equation}

\begin{proposition}[Weighted Neyman allocation]\label{prop:allocation}
\begin{enumerate}
\item[(a)] Over real $n_h>0$ with $\sum_hn_h=L$, $\Phi$ is minimised by
$n_h^\ast=L\,N_hA_h/\sum_kN_kA_k$, with $\Phi(n^\ast)=(\sum_hN_hA_h)^2/L-\sum_hN_hA_h^2$.
\item[(b)] With bounds $\ell_h\le n_h\le N_h$ ($\ell_h$ a minimum per stratum,
$\sum_h\ell_h\le L\le N$), the minimiser is $n_h=\min\{N_h,\max(\ell_h,N_hA_h/\lambda)\}$ with
$\lambda>0$ chosen so that $\sum_hn_h=L$.
\item[(c)] If an integer allocation $\tilde n$ satisfies $\tilde n_h\ge n_h-1$ and
$n_h\ge n_0>1$ for all $h$, then $\Phi_0(\tilde n)\le\{n_0/(n_0-1)\}\Phi_0(n)$, where
$\Phi_0(n)=\sum_hN_h^2A_h^2/n_h$ is the part of $\Phi$ that depends on $n$.
\item[(d)] Ignore finite-population corrections and compute variances with divisors $N_h$ and
$N$. With $w_t=1/\sigma^2_t$, $\sigma^2_t$ the population variance of target $t$'s influence
values,
\[
\Phi_0(n)\,L/N^2=\sum_t1/\mathrm{Eff}_t(n),\qquad \mathrm{Eff}_t=V^{SRS}_t/V_t .
\]
At the optimum the harmonic mean efficiency over the $|\mathcal T|$ targets equals
$|\mathcal T|/(\sum_hW_hA_h)^2\ge1$.
\end{enumerate}
\end{proposition}

Part (a) is Neyman's allocation \citep{neyman1934} for a weighted sum of variances; the proof is
one application of the Cauchy--Schwarz inequality,
$(\sum_hN_hA_h)^2=(\sum_h\{N_hA_hn_h^{-1/2}\}n_h^{1/2})^2\le\Phi_0(n)\,L$, with equality if and only
if $n_h\propto N_hA_h$. Part (b) follows from the Karush--Kuhn--Tucker conditions of a convex
problem; strata that would receive more than $N_h$ labels are enumerated and the remainder is
re-allocated, which is what the iterative capping in our implementation does. Part (c) bounds the
cost of rounding: with the minimum of 10 labels per stratum it is at most a factor $10/9$ on the
weighted variance, and much less in practice; exact integer optimisation is possible
\citep{wright2012} but unnecessary. Part (d) explains the default weights: with $w_t=1/S_t^2$ the
allocation minimises the sum of inverse efficiencies over all comparison estimands, so no target is
neglected because it happens to be measured on a small scale (the inequality in (d) is for the plain estimator, whose influence values define $\sigma^2_t$). Because the proposition holds for any
non-negative $S^2_{ht}$, it applies verbatim to the difference estimator with the residual variances.

In the design the $S^2_{ht}$ are unknown. They are replaced by their expectations under the
working model $Y_i\sim\text{Bernoulli}(\pi_i)$ independently, with the gradient evaluated at the
model-expected totals; up to terms of order $1/N_h$ this is the anticipated variance of
\citet{isaki1982}. By \eqref{eq:resid-moments} the anticipated stratum variances are
\begin{align}
\tilde A^{0\,2}_h&=\sum_tw_t\big(\bar\sigma^2_{ht}+S^2_{\tilde z,ht}\big)\quad\text{(plain)},\notag\\
\tilde A^{\psi\,2}_h&=\sum_tw_t\bar\sigma^2_{ht}\quad\text{(difference)},\qquad \bar\sigma^2_{ht}=N_h^{-1}\sum_{U_h}\pi_i(1-\pi_i)d_{it}^2.\label{eq:anticipated}
\end{align}
The plain rule is what the design uses (Section~\ref{sec:design}). The residual rule is the Neyman allocation for the prediction-powered residual
\citep{fisch2024stratppi}; it depends on the proxy only through $\pi_i(1-\pi_i)$ and the label-free $d_i$, so a proxy of any
sharpness gives a well-defined allocation.

\begin{proposition}[Allocation for the difference estimator]\label{prop:alloc-dif}
Let $n^\psi$ be the optimal allocation for the difference estimator ($A_h=A^\psi_h$, $A^{\psi\,2}_h=\sum_tw_tS^2_{e,ht}$) and $\tilde n^0$ the allocation obtained from the
plain anticipated variances. Put $\tilde\nu_h=\sum_tw_tS^2_{\tilde z,ht}/\sum_tw_t\bar\sigma^2_{ht}$, $x_h=(1+\tilde\nu_h)^{1/2}$ and assume the anticipated values are exact for the
residual variances ($\tilde A^\psi_h=A^\psi_h$). Ignoring caps, minima and rounding, the weighted variance of the difference estimator with the \emph{wrong} (plain) allocation satisfies
\begin{equation}\label{eq:alloc-loss}
\frac{\Phi_\psi(\tilde n^0)}{\Phi_\psi(n^\psi)}=\Big(\sum_h\omega_hx_h\Big)\Big(\sum_h\omega_h/x_h\Big)\le\frac{(x_{\max}+x_{\min})^2}{4x_{\max}x_{\min}},\qquad
\omega_h=\frac{N_hA^\psi_h}{\sum_kN_kA^\psi_k},
\end{equation}
with equality to one iff $\tilde\nu_h$ is constant over the strata with $\omega_h>0$. The same expressions hold, with $x_h$ replaced by $1/x_h$ and
$\omega_h\propto N_hA^0_h$, for the plain estimator allocated with the residual variances. Consequently, if $\tilde\nu_h\in[0,\nu_{\max}]$, the loss from
using the wrong rule is at most $\{1+(1+\nu_{\max})^{1/2}\}^2/\{4(1+\nu_{\max})^{1/2}\}$ in either direction, for instance $9/8$ for $\nu_{\max}=3$.
\end{proposition}

So the choice between the two rules is a second-order matter wherever the proxy's signal-to-noise ratio $\tilde\nu_h$ is similar across strata, which is the case for strata with
comparable disagreement structure, and it matters when $\tilde\nu_h$ is large in some strata and small in others (for example, very sharp proxies in the agreement strata of the
rate differences). Neither rule is optimal for the other estimator, and the plain rule remains the sensible default for the difference estimator: it is exactly Neyman for the plain estimator, is
robust to a proxy whose calibration is not trusted (the residual rule depends on $\pi_i(1-\pi_i)$ and gives zero weight to strata in which the proxy is
extreme), and its efficiency loss for the difference estimator is bounded by \eqref{eq:alloc-loss}. Section~\ref{sec:sim-dif} reports the predicted and the observed loss.

\subsection{Design effect and the disagreement rate}\label{sec:theory-deff}

Let $\sigma^2_{ht}=(1-1/N_h)S^2_{ht}$ and $\sigma^2_t=(1-1/N)S^2_t$ be the variances with divisor
$N_h$ and $N$. For a single target allocated by its own Neyman rule, and with $L/N\to0$, the
efficiency relative to SRS is
\begin{equation}\label{eq:deff-general}
\mathrm{Eff}_t=\frac{\sigma_t^2}{(\sum_hW_h\sigma_{ht})^2}\;\ge\;1,
\end{equation}
because
$\sigma_t^2=\sum_hW_h\sigma_{ht}^2+\sum_hW_h(\bar z_h-\bar z)^2\ge(\sum_hW_h\sigma_{ht})^2$; equality
holds only if all within-stratum standard deviations are equal and the stratum means coincide.
Combined with Proposition~\ref{prop:influence}, \eqref{eq:deff-general} gives explicit formulas for the plain estimator.
We state them for an idealised design in which, within each group, the strata separate the
instances on which $m$ and $j$ disagree (set $\mathcal D_g$, share $\rho_g$ of group $g$) from
those on which they agree (set $\mathcal A_g$). This is exactly the pattern design when $M=2$; for
$M>2$ it describes a design stratified for the pair $(m,j)$, and is discussed below.
Finite-population corrections are ignored ($L/N\to0$ and $n_h/N_h\to0$); they are retained in
the exact linearised efficiency used in Section~\ref{sec:simulation}.

\begin{corollary}[Design effect as a function of the disagreement rate]\label{prop:deff}
Let $\zeta_g=N_g/N$, $p_g=P_g/N_g$, $r_g=N_g^{-1}\sum_{G_i=g}Y_i|D_i|$ (positives on which the pair
disagrees, as a share of the group), $\pi_{\mathcal D,g}=r_g/\rho_g$,
$\pi_{\mathcal A,g}=(p_g-r_g)/(1-\rho_g)$ and $a_g=c_{\mathcal D,g}/\pi_{\mathcal D,g}$, where
$c_{\mathcal D,g}=\delta_gp_g/\rho_g$ is the mean of $Y_iD_i$ over $\mathcal D_g$; $a_g\in[-1,1]$
is the net share of disputed positives on which $m$ is right and $j$ wrong.
\begin{enumerate}
\item[(a)] For $\theta^{TPR}_{mj}$, with own-target Neyman allocation on the strata
$\{\mathcal A_g,\mathcal D_g\}_{g=0,1}$ and $L/N\to0$,
\begin{equation}\label{eq:deff-tpr}
\mathrm{Eff}^{TPR}=\frac{\sum_g(r_g-\delta_g^2p_g)/(\zeta_gp_g^2)}
{\Big[\sum_g\big\{\rho_g\sigma_{\mathcal D,g}+(1-\rho_g)|\delta_g|
\sqrt{\pi_{\mathcal A,g}(1-\pi_{\mathcal A,g})}\big\}/p_g\Big]^2},
\end{equation}
$\sigma^2_{\mathcal D,g}=\pi_{\mathcal D,g}-c_{\mathcal D,g}^2+\delta_g^2\pi_{\mathcal D,g}(1-\pi_{\mathcal D,g})
-2\delta_gc_{\mathcal D,g}(1-\pi_{\mathcal D,g})$. The same holds for $\theta^{FPR}_{mj}$ with $Y$
replaced by $1-Y$. For $\theta^{ERR}_{mj}$, with $\rho=\sum_g\zeta_g\rho_g$ and $\mu_g$ the mean of
$D_i(1-2Y_i)$ over $\mathcal D_g$,
\begin{equation}\label{eq:deff-err}
\mathrm{Eff}^{ERR}=\frac{\rho-(\theta^{ERR}_{mj})^2}{\big(\sum_g\zeta_g\rho_g\sqrt{1-\mu_g^2}\big)^2}.
\end{equation}
Any further split of $\mathcal A_g$ or $\mathcal D_g$ (for example $\mathcal A_g$ into $K=0$ and
$K=M$) can only increase \eqref{eq:deff-tpr} and \eqref{eq:deff-err}.
\item[(b)] \emph{Small targets.} If $\delta_g=0$ and $\rho_g=\rho$ for both groups, then
\begin{equation}\label{eq:deff-small}
\mathrm{Eff}^{TPR}=\frac1\rho\,\Gamma,\qquad
\Gamma=\frac{\sum_gr_g/(\zeta_gp_g^2)}{(\sum_g\sqrt{r_g}/p_g)^2}\ge1,
\end{equation}
where $\Gamma$ is the efficiency of Neyman allocation between the two groups alone; and if
$\theta^{ERR}_{mj}=0=\mu_g$, then $\mathrm{Eff}^{ERR}=1/\rho$.
\item[(c)] \emph{Target size.} For one group (drop $g$), with $\pi_{\mathcal D}$,
$\pi_{\mathcal A}$ and $a$ held fixed, as $\rho\to0$,
\begin{equation}\label{eq:deff-lambda}
\mathrm{Eff}^{TPR}=\frac{1+O(\rho)}{\rho\,(\beta+\lambda)^2},\qquad
\beta=\sqrt{1-a^2\pi_{\mathcal D}},\qquad
\lambda=\frac{|a|\sqrt{\pi_{\mathcal D}\,\pi_{\mathcal A}(1-\pi_{\mathcal A})}}{p}.
\end{equation}
Since $\delta=\rho\pi_{\mathcal D}a/p$, the target is proportional to $\rho$ and to $a$. A more
one-sided disagreement set (larger $|a|$, hence a larger target) changes the gain by the bounded
factor $(\beta+\lambda)^{-2}$, which is below one when the agreement region is noisy
($\lambda>1-\beta$), but it does not change the order $1/\rho$.
\item[(d)] \emph{When the gain vanishes.} As $\rho_g\to1$ the agreement strata disappear and
\eqref{eq:deff-tpr} and \eqref{eq:deff-err} tend to the efficiency of stratification by
group alone (for the rate differences with $\delta_g=0$ this is $\Gamma$; for the error difference
it equals 1 when $\mu_0=\mu_1$). The gain beyond group stratification is therefore governed by $1-\rho$: it is large when
the compared models disagree on a small share of the pool, it is reduced by a constant factor
when the pair's disagreements are one-sided ($|a_g|$ near 1) and the agreement region has
substantial label noise, and it disappears when the models disagree almost everywhere.
\end{enumerate}
\end{corollary}

The heuristic \emph{efficiency $\approx1/\rho$} is thus exact for an error difference with a
two-sided disagreement set ($\mu_g=0$), and correct up to the group factor $\Gamma$ and the
target-size factor $(\beta+\lambda)^{-2}$ for the rate differences. For a pair that disagrees on
5\% of the pool it promises a gain of the order of 20 relative to SRS, and of 2.5 for a pair that
disagrees on 40\%. These are gains of the stratification; the difference estimator multiplies them by the factor $G$ of
\eqref{eq:dif-deff}, which by \eqref{eq:dif-closed} is of order one to three and does not grow as $\rho\to0$: the difference estimator removes the part of the
label noise that a proxy of moderate sharpness can explain, whereas the stratification removes the unlabelled agreement region.

For $M>2$ the pattern strata $(G_i,K_i)$ do not separate the disagreements of a particular pair:
an instance with $K_i=4$ out of nine may or may not be disputed by $m$ and $j$. The general
formula \eqref{eq:deff-general} still applies with the stratum variances of
Proposition~\ref{prop:influence}. For an error difference whose disagreements are two-sided within
every stratum ($e_h=0$ for all $h$, hence $\theta^{ERR}_{mj}=0$), $\sigma^2_{ht}\propto\rho_h$
and \eqref{eq:deff-general} becomes
\begin{equation}\label{eq:deff-mixed}
\mathrm{Eff}^{ERR}=\frac{\rho}{(\sum_hW_h\sqrt{\rho_h})^2},\qquad 1\le\mathrm{Eff}^{ERR}\le\frac1\rho,
\end{equation}
since $\rho_h\le\sqrt{\rho_h}$ and $(\sum_hW_h\sqrt{\rho_h})^2\le\sum_hW_h\rho_h=\rho$. The bound $1/\rho$ is
attained only when every stratum is either entirely disputed or entirely agreed; the pattern
strata approach it to the extent that the pair's disagreements concentrate in a few values of
$K$. A lower bound of the same type holds for the pattern strata. Let $\rho_{\rm any}$ be the share
of instances on which not all $M$ models agree, i.e.\ with $0<K_i<M$. The pair can disagree only
there, so $\rho_h=0$ in the strata with $K\in\{0,M\}$, and Cauchy--Schwarz over the remaining strata
gives $(\sum_hW_h\sqrt{\rho_h})^2\le\rho_{\rm any}\,\rho$, hence
\begin{equation}\label{eq:deff-sandwich}
\frac{1}{\rho_{\rm any}}\;\le\;\mathrm{Eff}^{ERR}\;\le\;\frac{1}{\rho},
\end{equation}
with equality on the left when the pair's disagreements are spread evenly over the strata with
$0<K<M$ ($\rho_h=\rho/\rho_{\rm any}$ there). With nine models $\rho_{\rm any}$ is large (between
0.18 and 0.73 in our pools) while the pairwise $\rho$ is small, and the position of a pair within
this interval measures how well the counts $K$ locate its disagreements. Two further losses separate \eqref{eq:deff-tpr} and \eqref{eq:deff-err} from the efficiency of the implemented
design: a single allocation serves all $3\binom M2$ comparison estimands
(Proposition~\ref{prop:allocation}(d)), and it is computed from the proxy
(Section~\ref{sec:theory-proxy}). The exact linearised efficiency of the implemented design is
$N^2(1/L-1/N)S^2_t/V_t(n)$ with $V_t$ from \eqref{eq:linvar} (or \eqref{eq:resid} for the difference estimator); Section~\ref{sec:simulation}
compares it, and the closed forms above, with Monte Carlo results.

\subsection{Misspecified proxy}\label{sec:theory-proxy}

Let $\tilde n$ be the allocation computed from the anticipated values $\tilde A_h$, and $A_h$ the
values computed from the true labels, both with the same weights. This subsection concerns the proxy in the allocation; its role in the
difference estimator is Proposition~\ref{prop:dif-proxy}.

\begin{proposition}[Imperfect proxy]\label{prop:proxy}
\begin{enumerate}
\item[(a)] Proposition~\ref{prop:validity} holds for $\tilde n$ whatever the proxy, provided
$\tilde n_h\ge2$ in every stratum (guaranteed by the minimum of 10 labels).
\item[(b)] Ignoring caps, minima and rounding, the relative loss in the weighted variance is
\begin{equation}\label{eq:proxy-loss}
\frac{\Phi_0(\tilde n)}{\Phi_0(n^\ast)}=
\frac{(\sum_hN_h\tilde A_h)(\sum_hN_hA_h^2/\tilde A_h)}{(\sum_hN_hA_h)^2}
=\Big(\sum_h\omega_hx_h\Big)\Big(\sum_h\omega_h/x_h\Big)\ge1,
\end{equation}
with $\omega_h=N_hA_h/\sum_kN_kA_k$ and $x_h=\tilde A_h/A_h$. The loss is zero if and only if
$\tilde A_h\propto A_h$; a proxy that is wrong by a common factor costs nothing.
\item[(c)] If $x_{\min}\le x_h\le x_{\max}$, the loss is at most
$(x_{\max}+x_{\min})^2/(4x_{\max}x_{\min})$ (Kantorovich inequality). For instance, if every
anticipated value $\tilde A_h$ is within a factor 2 of $A_h$ up to a common scale
($x_{\max}/x_{\min}\le4$), the weighted variance is at most $25/16$ times the optimum on the same
strata, so the design keeps at least 64\% of the efficiency of the optimal allocation, whatever
the number of strata.
\item[(d)] Proportional allocation is the special case $\tilde A_h\equiv1$. Neither allocation
dominates the other for every population: \eqref{eq:proxy-loss} is smaller for the proxy
allocation if and only if
\[
\Big(\sum_h\omega_hx_h\Big)\Big(\sum_h\omega_h/x_h\Big)\le\Big(\sum_h\omega_h/A_h\Big)\Big(\sum_h\omega_hA_h\Big).
\]
A proxy that is badly wrong in a high-variance stratum can make the design worse than
proportional allocation, and, through the minimum of $\ell$ labels in that stratum, worse than SRS.
\item[(e)] A defensive mixture $n_h=\eta\,n^{\rm prop}_h+(1-\eta)\tilde n_h$, $0<\eta<1$, satisfies
$\Phi_0(n)\le\Phi_0(n^{\rm prop})/\eta$ and $\Phi_0(n)\le\Phi_0(\tilde n)/(1-\eta)$; since
proportional allocation is never worse than SRS up to the factor $(1-1/\min_hN_h)^{-1}$, the
mixture is never worse than $1/\eta$ times SRS in the weighted variance, up to the same factor.
\end{enumerate}
\end{proposition}

What can be guaranteed is therefore validity (part (a)) and a loss that is controlled by how far
the anticipated stratum standard deviations are from proportional to the true ones (parts (b) and
(c)). What cannot be guaranteed for the implemented rule is improvement over SRS for every
population; that requires a defensive component as in part (e). Two features make large losses
unlikely in our setting. The dominant part of the structure, namely which strata carry the
numerators of the comparison estimands, is determined by the predictions and does not depend on the
proxy (Proposition~\ref{prop:influence}); and the proxy enters $\tilde A_h$ only through
$\pi_h(1-\pi_h)$ and $r_h$, which are smooth in the proxy probabilities. Section~\ref{sec:simulation}
quantifies the loss for good, uninformative and adversarial proxies.
